\chapter{Decoders and Encoders}
\label{ch:coders}

We know how to build any combinational function from a table and simplify it on a
Karnaugh map. But for a security panel with four zones a practical
question arises: how do we turn on the lamp of a zone given its \emph{number}, and, the other way around,
how do we get the number of the zone that triggered and show it
on a digital display? These tasks come up so often that they have become
ready-made ``bricks'': the decoder, the encoder, and code converters. It pays
to recognize them at first sight. Let us start with the units that convert one
way of writing a number (a code) into another.

\begin{buildbox}
On the \emph{Learn} tab, in the book ``Selection and distribution,'' open the lesson
\emph{``Decoder''}: build a $2\to4$ decoder with an enable input.
Feed the input $a$ the values $0,1,2,3$ in turn and write in your notebook which bit
of the output $y$ equals one. Later in this chapter we will find out how this
``brick'' is built inside and where it is needed.
\end{buildbox}

\begin{figure}[H]
\centering
\begin{tikzpicture}[node distance=6mm and 32mm]
  \node[block, minimum width=34mm] (bin) {Binary code\\\footnotesize \code{10} (the number 2)};
  \node[oblock, minimum width=34mm, right=of bin] (oh) {One-hot code\\\footnotesize \code{0100} (a one in position 2)};
  \draw[arr] ([yshift=2mm]bin.east) -- node[above, font=\sffamily\small] {decoder} ([yshift=2mm]oh.west);
  \draw[arr] ([yshift=-2mm]oh.west) -- node[below, font=\sffamily\small] {encoder} ([yshift=-2mm]bin.east);
\end{tikzpicture}
\caption{The decoder and the encoder perform mutually inverse conversions}
\end{figure}

A \term{one-hot code} is a code in which, out of $N$
wires, exactly one carries a one, and the number of that wire is the
value being transmitted. It is convenient for \emph{selecting} something: turning on one of
several lamps, or one of several memory chips.

\section{The decoder}

A \term{decoder} has $n$ inputs and $2^n$ outputs. The inputs give
a binary number, and a 1 appears on the output with that number, while all the others
are 0. In schematics a decoder is labeled with the letters DC.

\begin{figure}[H]
\centering
\begin{minipage}{0.4\textwidth}
\centering
\begin{tikzpicture}
  \draw[thick, fill=FpgaLight] (0,0) rectangle (2.6,3.2);
  \draw[thick] (0.6,0) -- (0.6,3.2); \draw[thick] (2.0,0) -- (2.0,3.2);
  \node[font=\sffamily\bfseries\large] at (1.3,1.6) {DC};
  \foreach \i/\n in {0/0, 1/1} {
    \draw[wire] (-0.6,2.1-\i*0.8) node[left] {$a_\i$} -- (0,2.1-\i*0.8);
    \node[font=\sffamily\small] at (0.3,2.1-\i*0.8) {\pgfmathparse{int(2^\i)}\pgfmathresult};
  }
  \draw[wire] (-0.6,0.4) node[left] {$E$} -- (0,0.4); \node[font=\sffamily\small] at (0.3,0.4) {E};
  \foreach \i in {0,...,3} {
    \draw[wire] (2.6,2.6-\i*0.7) -- (3.2,2.6-\i*0.7) node[right] {$y_\i$};
    \node[font=\sffamily\small] at (2.3,2.6-\i*0.7) {\i};
  }
\end{tikzpicture}
\end{minipage}\hfill
\begin{minipage}{0.56\textwidth}
\centering
\begin{tabular}{ccc|cccc}
\toprule
$E$ & $a_1$ & $a_0$ & $y_0$ & $y_1$ & $y_2$ & $y_3$ \\
\midrule
0 & $\times$ & $\times$ & 0 & 0 & 0 & 0 \\
1 & 0 & 0 & \textbf{1} & 0 & 0 & 0 \\
1 & 0 & 1 & 0 & \textbf{1} & 0 & 0 \\
1 & 1 & 0 & 0 & 0 & \textbf{1} & 0 \\
1 & 1 & 1 & 0 & 0 & 0 & \textbf{1} \\
\bottomrule
\end{tabular}

\smallskip
{\small $\times$ means ``don't care, 0 or 1''}
\end{minipage}
\caption{A $2\to4$ decoder with an enable input $E$: symbol and truth table}
\label{fig:dc24}
\end{figure}

The input $E$ (Enable) lets you ``switch off'' the decoder: when $E=0$,
all outputs are zero, whatever number is applied.

How do we build a decoder from gates? Each output is a product
that equals one for exactly one combination of inputs:
\[
  y_0 = E\,\overline{a_1}\,\overline{a_0}, \quad
  y_1 = E\,\overline{a_1}\,a_0, \quad
  y_2 = E\,a_1\,\overline{a_0}, \quad
  y_3 = E\,a_1\,a_0.
\]

\begin{figure}[H]
\centering
\begin{tikzpicture}
  % inputs on the left; the buses a1, ~a1, a0, ~a0, E run downwards
  \draw[wire] (-0.4,0) node[left] {$a_1$} -- (1.0,0);
  \node[gNOT] (n1) at (2.2,0) {};
  \draw[wire] (1.0,0) -- (n1.input);
  \draw[wire] (n1.output) -- (3.4,0);
  \draw[wire] (-0.4,-1.0) node[left] {$a_0$} -- (4.6,-1.0);
  \node[gNOT] (n0) at (5.8,-1.0) {};
  \draw[wire] (4.6,-1.0) -- (n0.input);
  \draw[wire] (n0.output) -- (7.0,-1.0);
  \draw[wire] (-0.4,-2.0) node[left] {$E$} -- (8.2,-2.0);
  \foreach \x/\byy/\l in {1.0/0/a_1, 3.4/0/\overline{a_1}, 4.6/-1.0/a_0, 7.0/-1.0/\overline{a_0}, 8.2/-2.0/E} {
    \draw[very thick, FpgaBlue] (\x,\byy) -- (\x,-11.3);
    \node[below, font=\small] at (\x,-11.3) {$\l$};
  }
  \foreach \x/\byy in {1.0/0, 4.6/-1.0, 8.2/-2.0} \fill (\x,\byy) circle (1.3pt);
  % AND gates: y0 = ~a1 ~a0, y1 = ~a1 a0, y2 = a1 ~a0, y3 = a1 a0
  \foreach \i/\xp/\xq in {0/3.4/7.0, 1/3.4/4.6, 2/1.0/7.0, 3/1.0/4.6} {
    \pgfmathsetmacro{\y}{-3.8-\i*2.0}
    \node[gAND, logic gate inputs=nnn, minimum height=13.5mm] (g\i) at (11.0,\y) {}; % fig-ok: 3-input gate, 13.5 mm
    \draw[wire] (g\i.input 1) -- (\xp,0 |- g\i.input 1) node[circle, fill, inner sep=1.1pt] {};
    \draw[wire] (g\i.input 2) -- (\xq,0 |- g\i.input 2) node[circle, fill, inner sep=1.1pt] {};
    \draw[wire] (g\i.input 3) -- (8.2,0 |- g\i.input 3) node[circle, fill, inner sep=1.1pt] {};
    \draw[wire] (g\i.output) -- ++(0.6,0) node[right] {$y_\i$};
  }
\end{tikzpicture}
\caption{A $2\to4$ decoder built from AND and NOT gates}
\label{fig:dc24_gates}
\end{figure}

\begin{figure}[H]
\centering
\begin{tikztimingtable}[timing/slope=0.05, xscale=2.2, yscale=1.4, timing/font=\sffamily\small, timing/rowdist=3.2ex]
  $a_1a_0$ & 2D{00} 2D{01} 2D{10} 2D{11} 2D{00} \\
  $y_0$ & 2H 6L 2H \\
  $y_1$ & 2L 2H 6L \\
  $y_2$ & 4L 2H 4L \\
  $y_3$ & 6L 2H 2L \\
\end{tikztimingtable}
\caption{As the numbers are swept, the one ``moves'' from one output to another}
\end{figure}

\subsection{Where decoders are used}

\begin{itemize}
  \item \textbf{Selecting a device by address.} In a computer, many memory and
        peripheral chips are connected to one bus. A decoder uses the
        high bits of the address to select which of them should respond right now.
  \item \textbf{Driving displays}: turning on one of several
        LEDs, display digits, or matrix rows.
  \item \textbf{Selecting an instruction in a processor}: the operation code is decoded into
        the control signals ``add,'' ``read from memory,'' and so on.
\end{itemize}

\begin{figure}[H]
\centering
\begin{tikzpicture}[node distance=5mm]
  \node[hblock, minimum height=30mm, minimum width=18mm] (cpu) {Processor};
  \node[block, right=22mm of cpu, minimum height=26mm] (dc) {DC\\2$\to$4};
  \foreach \i/\t in {0/ROM, 1/RAM, 2/Input port, 3/Output port} {
    \node[oblock, minimum width=24mm, minimum height=7mm] (d\i) at (9.0,1.65-\i*1.1) {\t};
  }
  \draw[arr] (cpu.east) -- node[above, lbl] {A[15:14]} (dc.west);
  \foreach \i in {0,...,3} \draw[arr] ($(dc.north east)!{(\i+0.5)/4}!(dc.south east)$) -- ++(0.6,0) |- (d\i.west);
  \foreach \i in {0,...,3} \node[lbl, right=1mm of d\i.east, anchor=west] {CS$_\i$};
\end{tikzpicture}
\caption{An address decoder: the two high bits of the address select one of four chips}
\end{figure}

\subsection{Cascading: a large decoder from small ones}

The enable input $E$ is needed not only to ``switch off'' a decoder: with its help,
large decoders are assembled from small ones. A $3\to8$ decoder is made from
two $2\to4$ decoders: the low two bits of the number $a_1a_0$ go to
both, and the high bit $a_2$ decides which of the two ``wakes up.'' The first
is enabled by $\overline{a_2}$, the second by $a_2$. Then, when $a_2=0$, the
first works and lights one of the lines $y_0\ldots y_3$, and when $a_2=1$ the second works and
lights one of $y_4\ldots y_7$. The same technique (the first level selects
a chip, the second selects a line) is used when selecting a device by address.

\begin{figure}[H]
\centering
\begin{tikzpicture}
  \node[block, minimum width=22mm, minimum height=16mm] (d0) at (5,1.3) {DC\\2$\to$4};
  \node[block, minimum width=22mm, minimum height=16mm] (d1) at (5,-1.5) {DC\\2$\to$4};
  \draw[arr] (0,0.1) node[left] {$a_1a_0$} -- (2.4,0.1);
  \fill (2.4,0.1) circle (1.3pt);
  \draw[arr] (2.4,0.1) |- ($(d0.west)+(0,-0.3)$);
  \draw[arr] (2.4,0.1) |- ($(d1.west)+(0,0.3)$);
  \draw[arr] (0,2.9) node[left] {$\overline{a_2}$} -- (5,2.9) -- node[right, lbl] {$E$} (d0.north);
  \draw[arr] (0,-3.3) node[left] {$a_2$} -- (5,-3.3) -- node[right, lbl] {$E$} (d1.south);
  \draw[arr] ($(d0.east)+(0,0)$) -- ++(1.4,0) node[right] {$y_0\ldots y_3$};
  \draw[arr] ($(d1.east)+(0,0)$) -- ++(1.4,0) node[right] {$y_4\ldots y_7$};
\end{tikzpicture}
\caption{A $3\to8$ decoder from two $2\to4$ decoders with an enable input}
\label{fig:dc38}
\end{figure}

\begin{diybox}[A decoder in Logra]
\textbf{Task.} \stars{1} Build a $2\to4$ decoder and check all
eight combinations $(E,a)$. \stars{2} Compare it with your own circuit of
AND and NOT gates (Figure~\ref{fig:dc24_gates}): the number of gates and the
result. \stars{3} Build a $3\to8$ decoder from two $2\to4$ decoders as in
Figure~\ref{fig:dc38}; acceptance criterion: for each of the eight numbers
exactly one line is lit, and for $E=0$ none.

\textbf{In Logra.} \stars{1} The lesson \emph{``Decoder''} (book ``Selection and
distribution''); the example \code{decoder}. \stars{2} The \emph{Designer} tab:
AND and NOT blocks. \stars{3} The \emph{Designer} tab: two ``Decoder'' blocks with
enable switched on, a ``Bus slice'' block to extract $a_1a_0$ and $a_2$
(lesson ``Bus: a slice of bits''), and a NOT gate.

\textbf{Expected result.} For the number 5 ($a=101$), $y_5$ lights up and
the other outputs are zeros; for 2 ($a=010$), $y_2$ lights up. The example
\code{binary-onehot} shows that ``extra'' numbers give an empty
combination if you do not provide for them.

\textbf{How to check.} The \emph{Waveforms} tab: drive $a$ with a counter
$0$--$7$; the one should ``run'' along the outputs. The \emph{RTL} tab: compare
your circuit with a single ``Decoder'' block.
\end{diybox}

\section{The seven-segment display}

A seven-segment display consists of seven LED strips, labeled
with the letters $a$ to $g$. By lighting different segments, you can show any digit.
The unit that converts the binary code of a digit into signals for the segments is also
called a decoder (although it does not produce a one-hot code).

\newcommand{\sevenseg}[8]{%
  \begin{tikzpicture}[scale=0.32, seg/.style={line width=2.6pt, line cap=round}]
    \def\segc##1{\ifnum##1=1 \def\cc{FpgaRed}\else\def\cc{gray!20}\fi}
    \segc{#1}\draw[seg, color=\cc] (0.2,4) -- (1.8,4);   % a
    \segc{#2}\draw[seg, color=\cc] (2,3.8) -- (2,2.2);   % b
    \segc{#3}\draw[seg, color=\cc] (2,1.8) -- (2,0.2);   % c
    \segc{#4}\draw[seg, color=\cc] (0.2,0) -- (1.8,0);   % d
    \segc{#5}\draw[seg, color=\cc] (0,1.8) -- (0,0.2);   % e
    \segc{#6}\draw[seg, color=\cc] (0,3.8) -- (0,2.2);   % f
    \segc{#7}\draw[seg, color=\cc] (0.2,2) -- (1.8,2);   % g
    \node[font=\sffamily\footnotesize] at (1,-1.1) {#8};
  \end{tikzpicture}}

\begin{figure}[H]
\centering
\begin{tikzpicture}[scale=0.8, seg/.style={line width=4pt, line cap=round, FpgaRed}]
  \draw[seg] (0.25,4) -- (1.75,4) node[midway, above=1mm, text=FpgaDark, font=\sffamily\small] {a};
  \draw[seg] (2,3.75) -- (2,2.25) node[midway, right=1mm, text=FpgaDark, font=\sffamily\small] {b};
  \draw[seg] (2,1.75) -- (2,0.25) node[midway, right=1mm, text=FpgaDark, font=\sffamily\small] {c};
  \draw[seg] (0.25,0) -- (1.75,0) node[midway, below=1mm, text=FpgaDark, font=\sffamily\small] {d};
  \draw[seg] (0,1.75) -- (0,0.25) node[midway, left=1mm, text=FpgaDark, font=\sffamily\small] {e};
  \draw[seg] (0,3.75) -- (0,2.25) node[midway, left=1mm, text=FpgaDark, font=\sffamily\small] {f};
  \draw[seg] (0.25,2) -- (1.75,2) node[midway, above=0.5mm, text=FpgaDark, font=\sffamily\small] {g};
\end{tikzpicture}
\hspace{12mm}
\sevenseg1111110{0}\,\sevenseg0110000{1}\,\sevenseg1101101{2}\,\sevenseg1111001{3}\,\sevenseg0110011{4}\,%
\sevenseg1011011{5}\,\sevenseg1011111{6}\,\sevenseg1110000{7}\,\sevenseg1111111{8}\,\sevenseg1111011{9}
\caption{A seven-segment display: segment labels and the shapes of the digits}
\end{figure}

\begin{table}[H]
\centering
\caption{Truth table of the seven-segment display decoder (1 means the segment is lit)}
\label{tab:7seg}
\small
\begin{tabular}{c|cccc|ccccccc}
\toprule
Digit & $x_3$ & $x_2$ & $x_1$ & $x_0$ & a & b & c & d & e & f & g \\
\midrule
0 & 0&0&0&0 & 1&1&1&1&1&1&0 \\
1 & 0&0&0&1 & 0&1&1&0&0&0&0 \\
2 & 0&0&1&0 & 1&1&0&1&1&0&1 \\
3 & 0&0&1&1 & 1&1&1&1&0&0&1 \\
4 & 0&1&0&0 & 0&1&1&0&0&1&1 \\
5 & 0&1&0&1 & 1&0&1&1&0&1&1 \\
6 & 0&1&1&0 & 1&0&1&1&1&1&1 \\
7 & 0&1&1&1 & 1&1&1&0&0&0&0 \\
8 & 1&0&0&0 & 1&1&1&1&1&1&1 \\
9 & 1&0&0&1 & 1&1&1&1&0&1&1 \\
\bottomrule
\end{tabular}
\end{table}

Each column a--g is a separate logic function of four
variables. It can be simplified on a Karnaugh map and built from gates, as in
Chapters~\ref{ch:comb} and~\ref{ch:minimize}. The codes $1010$--$1111$ are
not digits, so nothing is required of them: for the
segment functions they are don't-care states, and they make the
minimal forms noticeably shorter.

\textbf{Example: segment $e$.} In Table~\ref{tab:7seg}, segment $e$ is lit on
the digits $0,2,6,8$. Let us draw the map for the function $e(x_3,x_2,x_1,x_0)$,
marking the don't-care cells $10$--$15$ with the symbol $\times$.

\begin{figure}[H]
\centering
\begin{tikzpicture}[cell/.style={draw, minimum size=10mm, font=\sffamily\large}]
  \foreach \c/\lab in {0/00, 1/01, 2/11, 3/10} \node[font=\sffamily\small] at (\c,0.85) {\lab};
  \foreach \r/\lab in {0/00, 1/01, 2/11, 3/10} \node[font=\sffamily\small] at (-0.95,-\r) {\lab};
  \node[font=\sffamily\small, anchor=south east] at (-0.5,0.9) {$x_3x_2\backslash x_1x_0$};  % e: 0,2,6,8 = 1; 1,3,4,5,7,9 = 0; 10..15 = x
  \foreach \r/\c/\v in {0/0/1, 0/1/0, 0/2/0, 0/3/1,
                        1/0/0, 1/1/0, 1/2/0, 1/3/1,
                        2/0/x, 2/1/x, 2/2/x, 2/3/x,
                        3/0/1, 3/1/0, 3/2/x, 3/3/x}
    {\def\cc{white}
     \if 1\v \def\cc{FpgaLight}\fi
     \if x\v \def\cc{gray!25}\fi
     \node[cell, fill=\cc] at (\c,-\r) {\if x\v$\times$\else\v\fi};}
  \foreach \x/\y in {0/0, 3/0, 0/-3, 3/-3}
    \draw[very thick, FpgaRed, rounded corners=4pt] (\x-0.44,\y+0.44) rectangle (\x+0.44,\y-0.44);
  \draw[very thick, FpgaGreen, rounded corners=4pt] (2.62,0.38) rectangle (3.38,-3.38);
  \node[font=\sffamily\small, text=FpgaRed, anchor=west] at (4.4,0.2) {$\overline{x_2}\,\overline{x_0}$ --- four corners};
  \node[font=\sffamily\small, text=FpgaGreen, anchor=west] at (4.4,-0.6) {$x_1\overline{x_0}$ --- right column};
  \node[font=\sffamily\small\bfseries, anchor=west] at (4.4,-1.6) {$e=\overline{x_0}\,(\overline{x_2}+x_1)$};
\end{tikzpicture}
\caption{Segment $e$ on a Karnaugh map: the don't-care cells allow two large groups of four cells}
\label{fig:seg_e}
\end{figure}

The four corners (cells $0,2,8,10$) give $\overline{x_2}\,\overline{x_0}$,
and the right column (cells $2,6,10,14$) gives $x_1\overline{x_0}$. We add them and
factor out the common term: $e=\overline{x_0}(\overline{x_2}+x_1)$, just three
gates instead of seven terms of four letters each. Let us check on digits:
for $4$ ($0100$) we get $1\cdot(0+0)=0$, and segment $e$ is off; for $6$
($0110$) we get $1\cdot(0+1)=1$, and it is lit. Similarly, segment $c$ is off only
on the digit 2, and its minimal form is $c=x_2+\overline{x_1}+x_0$.

\begin{warnbox}[Active level]
Many real displays have a common anode: a segment lights not when the
decoder output is one but when it is \emph{zero}. Then all seven outputs must be
inverted. If on the board you saw a ``negative'' instead of a digit, the likely
cause is exactly this, not a mistake in the table.
\end{warnbox}

\begin{diybox}[A digit on seven segments]
\textbf{Task.} \stars{1} Show the digits $0\ldots9$ on the display.
\stars{2} Add the hexadecimal digits $A\ldots F$ and see
how they differ from the ``digits.'' \stars{3} Find the minimal form of
segment $c$ and build it from two gates: $c=x_2+\overline{x_1}+x_0$.

\textbf{In Logra.} \stars{1} The lesson \emph{``A digit on seven segments''} (book
``Displays''); the example \code{segment-decoder} (the example library).
\stars{2} The lesson \emph{``Digits from 0 to F''}. \stars{3} The \emph{Designer}
tab: an OR block with three inputs and an inverter.

\textbf{Expected result.} A seven-bit output; the least significant bit is segment
$a$, the most significant is $g$. For the digit 1 you get \code{0000110} ($b$ and $c$ lit).
For the codes $10$--$15$ the ``BCD\,$\to$\,7 segments'' block turns all
segments off, that is, it chooses zeros for the don't-care states.

\textbf{How to check.} The \emph{Waveforms} tab: sweep $0$--$9$
and compare with Table~\ref{tab:7seg}. For your own segment
$c$, compare the output with one column of the table on all ten digits; on
$10$--$15$ the result does not matter.
\end{diybox}

\section{The encoder}

An \term{encoder} (labeled CD in schematics) performs the inverse
conversion: from a one-hot code on $2^n$ inputs it produces the $n$-bit
binary number of the line that carries the one.

\begin{figure}[H]
\centering
\begin{minipage}{0.38\textwidth}
\centering
\begin{tikzpicture}
  \draw[thick, fill=FpgaLight] (0,0) rectangle (2.6,3.0);
  \draw[thick] (0.6,0) -- (0.6,3.0); \draw[thick] (2.0,0) -- (2.0,3.0);
  \node[font=\sffamily\bfseries\large] at (1.3,1.5) {CD};
  \foreach \i in {0,...,3} {
    \draw[wire] (-0.6,2.4-\i*0.65) node[left] {$x_\i$} -- (0,2.4-\i*0.65);
    \node[font=\sffamily\small] at (0.3,2.4-\i*0.65) {\i};
  }
  \foreach \i/\w in {0/1, 1/2} {
    \draw[wire] (2.6,1.9-\i*0.9) -- (3.2,1.9-\i*0.9) node[right] {$y_\i$};
    \node[font=\sffamily\small] at (2.3,1.9-\i*0.9) {\w};
  }
\end{tikzpicture}
\end{minipage}\hfill
\begin{minipage}{0.58\textwidth}
\centering
\begin{tabular}{cccc|cc}
\toprule
$x_0$ & $x_1$ & $x_2$ & $x_3$ & $y_1$ & $y_0$ \\
\midrule
1 & 0 & 0 & 0 & 0 & 0 \\
0 & 1 & 0 & 0 & 0 & 1 \\
0 & 0 & 1 & 0 & 1 & 0 \\
0 & 0 & 0 & 1 & 1 & 1 \\
\bottomrule
\end{tabular}
\end{minipage}
\caption{A $4\to2$ encoder: symbol and truth table}
\end{figure}

The table shows that $y_0$ equals 1 when line 1 or 3 is active; $y_1$ when
line 2 or 3 is active. So the encoder is built from OR gates alone:
\[
  y_0 = x_1 + x_3, \qquad y_1 = x_2 + x_3.
\]

\begin{figure}[H]
\centering
\begin{tikzpicture}
  \foreach \i in {0,...,3} \draw[very thick, FpgaBlue] (\i*0.7,0.6) node[above, text=FpgaDark] {$x_\i$} -- (\i*0.7,-2.4);
  \node[gOR] (o0) at (4.2,-0.4) {};
  \node[gOR] (o1) at (4.2,-2.0) {};
  \draw[wire] (o0.input 1) -- (0.7,0 |- o0.input 1) node[circle, fill, inner sep=1.1pt] {};
  \draw[wire] (o0.input 2) -- (2.1,0 |- o0.input 2) node[circle, fill, inner sep=1.1pt] {};
  \draw[wire] (o1.input 1) -- (1.4,0 |- o1.input 1) node[circle, fill, inner sep=1.1pt] {};
  \draw[wire] (o1.input 2) -- (2.1,0 |- o1.input 2) node[circle, fill, inner sep=1.1pt] {};
  \draw[wire] (o0.output) -- ++(0.6,0) node[right] {$y_0$};
  \draw[wire] (o1.output) -- ++(0.6,0) node[right] {$y_1$};
  \node[lbl, text=FpgaGray] at (0,-3.1) {unused};
\end{tikzpicture}
\caption{A $4\to2$ encoder from two OR gates}
\end{figure}

\subsection{The priority encoder}

What happens if a one appears on two inputs at once, for example if
buttons 1 and 2 are pressed at the same time? The simple encoder outputs $y_1y_0 = 11$, that is,
the number 3, which is plainly an error. So in practice a
\term{priority encoder} is used: it outputs the number of the \emph{highest} active
input and ignores the others. An additional output $V$ (valid) shows
that at least one input is active.

\begin{table}[H]
\centering
\caption{A $4\to2$ priority encoder ($\times$ means the value does not matter)}
\begin{tabular}{cccc|ccc}
\toprule
$x_3$ & $x_2$ & $x_1$ & $x_0$ & $y_1$ & $y_0$ & $V$ \\
\midrule
0 & 0 & 0 & 0 & 0 & 0 & 0 \\
0 & 0 & 0 & 1 & 0 & 0 & 1 \\
0 & 0 & 1 & $\times$ & 0 & 1 & 1 \\
0 & 1 & $\times$ & $\times$ & 1 & 0 & 1 \\
1 & $\times$ & $\times$ & $\times$ & 1 & 1 & 1 \\
\bottomrule
\end{tabular}
\end{table}

The expressions: $y_1 = x_3 + x_2$, \ $y_0 = x_3 + \overline{x_2}\,x_1$, \
$V = x_3 + x_2 + x_1 + x_0$.

\begin{infobox}[Where encoders are used]
\begin{itemize}
  \item \textbf{Keyboards}: the number of the pressed key is turned into a code.
  \item \textbf{Interrupt controllers}: of several simultaneous requests,
        the number of the most important one is passed to the processor.
  \item \textbf{Position sensors}: from the readings of a row of photosensors
        you get a number, the coordinate.
\end{itemize}
\end{infobox}

Why the highest and not the lowest? The convention is arbitrary; what matters
is only that it is \emph{the same} everywhere: in the keyboard, in the interrupt
controller, in the security panel. In the panel this means: if zones~1 and~3 trigger
at the same time, the display will show the number~3, and the guard
will go where it is ``more important.'' The output $V$ is needed to tell ``zone~0
triggered'' from ``nothing triggered'': in both cases $y_1y_0=00$.

\begin{diybox}[Encoder and priority]
\textbf{Task.} \stars{1} For the inputs $x_3x_2x_1x_0=0101$ and $1010$,
determine the output of the priority encoder ($y_1y_0$, $V$). \stars{2}
Build an encoder from OR gates using the formulas $y_0=x_1+x_3$,
$y_1=x_2+x_3$ and compare it with the priority encoder when there are two ones on the input.
\stars{3} Add the output $V$ to the priority encoder using gates
(the formula $V=x_3+x_2+x_1+x_0$).

\textbf{In Logra.} \stars{1} The lesson \emph{``Priority encoder''} (book
``Selection and distribution''); the example \code{priority} (the example library).
\stars{2}, \stars{3} The \emph{Designer} tab: OR, NOT, and AND blocks.

\textbf{Expected result.} For $0101$: $y_1y_0=10$, $V=1$; for $1010$:
$y_1y_0=11$, $V=1$. With two ones the simple encoder outputs
the \emph{wrong} number: for example, with $x_1=x_2=1$ it gives $11$, that is,
the number~3, which is not on the input.

\textbf{How to check.} The \emph{Waveforms} tab: apply all
16 combinations to $x$; make sure that with a one in $x_3$ you always get $y=11$,
and that $V=0$ only when $x=0000$.
\end{diybox}

\section{Code converters}

Decoders and encoders are special cases of \term{code converters}.
Another useful example is the conversion of binary code to \term{Gray code}, in
which adjacent numbers differ in exactly one bit. This matters for
shaft angle sensors: on the transition from 3 (\code{011}) to 4 (\code{100}), three bits
change at once in binary code, and at the intermediate moment the sensor may output
anything. In Gray code only one bit changes. (This is the same
mechanism as the glitches in Chapter~\ref{ch:comb}: several signals change
at different times, and a false value flashes between them.)

\begin{figure}[H]
\centering
\begin{minipage}{0.45\textwidth}
\centering
\begin{tabular}{c c c}
\toprule
\textbf{Number} & \textbf{Binary} & \textbf{Gray} \\
\midrule
0 & \code{000} & \code{000} \\
1 & \code{001} & \code{001} \\
2 & \code{010} & \code{011} \\
3 & \code{011} & \code{010} \\
4 & \code{100} & \code{110} \\
5 & \code{101} & \code{111} \\
6 & \code{110} & \code{101} \\
7 & \code{111} & \code{100} \\
\bottomrule
\end{tabular}
\end{minipage}\hfill
\begin{minipage}{0.5\textwidth}
\centering
\begin{tikzpicture}
  \draw[very thick, FpgaBlue] (0,0) node[left, text=FpgaDark] {$b_2$} -- (3.4,0) node[right, text=FpgaDark] {$g_2$};
  \node[gXOR] (x1) at (2.4,-1.0) {};
  \node[gXOR] (x0) at (2.4,-2.5) {};
  \fill (0.6,0) circle (1.3pt);
  \draw[wire] (x1.input 1) -- (0.6,0 |- x1.input 1) -- (0.6,0);
  \draw[wire] (x1.input 2) -- (0,0 |- x1.input 2) node[left] {$b_1$};
  \fill (1.0,0 |- x1.input 2) circle (1.3pt);
  \draw[wire] (x0.input 1) -- (1.0,0 |- x0.input 1) -- (1.0,0 |- x1.input 2);
  \draw[wire] (x0.input 2) -- (0,0 |- x0.input 2) node[left] {$b_0$};
  \draw[wire] (x1.output) -- ++(0.6,0) node[right] {$g_1$};
  \draw[wire] (x0.output) -- ++(0.6,0) node[right] {$g_0$};
\end{tikzpicture}

\small $g_i = b_{i+1} \oplus b_i$
\end{minipage}
\caption{Binary code and Gray code; a converter built from ``exclusive OR'' gates}
\end{figure}

The inverse conversion goes as a ``chain'': the high bit stays, and each
next bit equals the XOR of the binary bit already obtained and the current Gray bit,
$b_{n-1}=g_{n-1}$, $b_i=b_{i+1}\oplus g_i$. For example, the Gray code \code{110}
gives $b_2=1$, $b_1=1\oplus1=0$, $b_0=0\oplus0=0$, that is, the binary
\code{100} (the number~4). The forward conversion is simpler: each Gray
bit depends only on two adjacent binary bits, so all the XORs work
in parallel, with no chain.

\begin{diybox}[Gray code]
\textbf{Task.} \stars{1} Write the numbers $0\ldots7$ in Gray code and
check that adjacent ones differ in one bit. \stars{2} Build a
converter from binary code to Gray code from XOR gates for three bits.
\stars{3} Build the inverse converter and make sure that the double
conversion returns the original number.

\textbf{In Logra.} \stars{1} The lesson \emph{``Gray code with a block''} (book ``Comparison
and codes''). \stars{2} The \emph{Designer} tab: two XOR blocks. The example
\code{gray} (the example library): compare \code{binary} and \code{gray} on
the transition $3\to4$. \stars{3} The example \code{binary-onehot} as a model of
a ``code $\to$ code'' conversion.

\textbf{Expected result.} On the transition $3\to4$, three
bits of the binary code change and one bit of the Gray code. The double conversion
returns the number.

\textbf{How to check.} The \emph{Waveforms} tab: drive the input $b$ with a
counter $0$--$7$ and count how many Gray bits change on
each step; there should be exactly one change everywhere.
\end{diybox}

\section{The same blocks in Verilog}
\label{sec:coders-verilog}

All the blocks of this chapter are convenient to describe not with a drawing
but with text (the Verilog language is covered in
Chapter~\ref{ch:verilog}). The reading rules are short.
\code{assign} means ``the output is always equal to this expression''.
An \code{always @*} block with \code{case} or \code{if} reads as ``a table or
a chain of conditions that selects one value''; inside it we write the
blocking assignment \code{=} and \emph{assign a value to every output in every
branch}.

How to run any example of this section in Logra: the \emph{Code} tab, paste
the module (or open a file from \code{code/b4\_coders}), press
\emph{Synthesize} (Ctrl+Enter); the \emph{RTL} tab shows the resulting
circuit, the \emph{Waveforms} tab sweeps the inputs. Next to each module there
is a testbench \code{tb\_*.v}: it sweeps the inputs, compares the outputs with a
reference, and prints \code{OK}.

\subsection{The $4\to2$ encoder with case}

The encoder table turns into a \code{case} almost literally: the one-hot
code on the left, the line number on the right.

\vfile{\codepath/b4_coders/enc4to2.v}{code/b4\_coders/enc4to2.v}

\textbf{How to read it.} \code{case (x)} compares the input with four
patterns from top to bottom; the first match sets \code{y}. The \code{default}
branch takes everything else: zero ones and several ones. Without it, for
those codes \code{y} would stay ``as it was'', that is, memory would appear
(see the subsection on mistakes below). Compare with the ``gates'' version
we drew at the beginning of the section:

\vfile{\codepath/b4_coders/enc4to2_or.v}{code/b4\_coders/enc4to2\_or.v}

On one-hot codes the two versions are identical. For $x=0110$ the OR version
gives $y=11$ (a wrong number, as we predicted), while the \code{case}
version honestly ``does not know'' and gives $00$: the behaviour on
forbidden inputs is now written down explicitly instead of coming out by accident.

\vfile{\codepath/b4_coders/tb_enc4to2.v}{code/b4\_coders/tb\_enc4to2.v}

\textbf{Scenario on the Waveforms tab.} Set $x$ to
\code{0001, 0010, 0100, 1000}: $y$ must go through $0,1,2,3$ in both modules;
then $x=0110$, and find how the two $y$ differ.

\subsection{The priority encoder: if/else if and casez}

Priority in text is the \emph{order of checks}. The first check that turns out
true wins, and the rest are not even looked at.

\vfile{\codepath/b4_coders/prio4_if.v}{code/b4\_coders/prio4\_if.v}

\textbf{How to read it.} First \code{v = 1} by default. Then the chain:
``if $x_3$ is active, number 3; otherwise, if $x_2$, number 2; \ldots;
otherwise (nothing there) number 0 and $v=0$.'' Every branch assigns both \code{y}
and \code{v}: otherwise a latch would appear. The same function with
\code{casez}, where the sign \code{?} means ``anything may be here'', which
is exactly the $\times$ of the priority encoder table:

\vfile{\codepath/b4_coders/prio4_casez.v}{code/b4\_coders/prio4\_casez.v}

The \code{casez} rows are checked from top to bottom, like an \code{if}
chain: for $x=0110$ the first row does not match (the top bit is zero), the
second one, \code{01??}, matches, and $y=2$. The two descriptions are
equivalent; \code{casez} repeats the table more clearly, \code{if} is better
when the conditions are more complex than a ``mask''.

\vfile{\codepath/b4_coders/tb_prio4.v}{code/b4\_coders/tb\_prio4.v}

\textbf{Scenario on the Waveforms tab.} Drive $x$ with a counter $0$--$15$:
for $x=0$ the output $v=0$; for any one in $x_3$ always $y=3$; for
$x=0110$ you get $y=2$, and both versions of the module give identical plots.

\subsection{Decoders: assign and case}

A $2\to4$ decoder can be written with one \code{assign} per output
(the formulas from the beginning of the chapter) or as a \code{case} table:

\vfile{\codepath/b4_coders/dec2to4_assign.v}{code/b4\_coders/dec2to4\_assign.v}

\vfile{\codepath/b4_coders/dec2to4_case.v}{code/b4\_coders/dec2to4\_case.v}

\textbf{How to read it.} In the first module each line is a product
$E\,\overline{a_1}\,\overline{a_0}$ and so on, the sign \code{\textasciitilde} is NOT.
In the second the enable is checked first: for \code{en = 0} all outputs are
zero, otherwise \code{case (a)} picks a table row. Next is a $3\to8$ decoder
in one line: a one is shifted left by \code{a} positions, and for \code{en = 0}
the output is zero; and the same decoder assembled from two $2\to4$ blocks as
in Figure~\ref{fig:dc38}.

\vfile{\codepath/b4_coders/dec3to8.v}{code/b4\_coders/dec3to8.v}

\vfile{\codepath/b4_coders/dec3to8_cascade.v}{code/b4\_coders/dec3to8\_cascade.v}

In the cascade version the low two bits of the number go to both blocks and
the top bit enables one of them: \code{en \& \textasciitilde a[2]} turns on the
lower half ($y_0\ldots y_3$), \code{en \& a[2]} the upper one. This is exactly
the circuit of Figure~\ref{fig:dc38}.

\textbf{Scenario.} The testbench \code{tb\_dec.v} compares all four
descriptions with the reference on all 16 combinations $(E,a)$. On the
\emph{Waveforms} tab drive $a$ with a counter $0$--$7$ with $en=1$: the one
``runs'' along the outputs of \code{dec3to8}; for the cascade the \emph{RTL}
tab shows two identical blocks and the \code{\&}, \code{\textasciitilde} gates.

\subsection{The seven-segment decoder}

Table~\ref{tab:7seg} turns into a \code{case} literally row by row; the seven
outputs are collected into one bus \code{seg = \{g,f,e,d,c,b,a\}}
(the least significant bit is segment $a$).

\vfile{\codepath/b4_coders/seg7.v}{code/b4\_coders/seg7.v}

\textbf{How to read it.} For the digit 1 we get \code{7'b0000110}, that is,
$b$ and $c$ are lit, as in the table. The \code{default} branch covers the
codes $10$--$15$ and blanks the display. This is the very solution for
don't-care states we discussed above; instead of zeros one could write \code{7'bx},
and the optimizer would choose the values itself (see Chapter~\ref{ch:minimize}).
The testbench \code{tb\_seg7.v} checks all 16 codes.

\textbf{Scenario.} The \emph{Waveforms} tab: drive $x$ with a counter $0$--$15$;
on the digits $0$--$9$ compare the output with the table, on $10$--$15$ it is zero.
The \emph{RTL} tab: see what this \code{case} turned into: the digits are
not ``taken apart'' into seven separate circuits, it is one table.

\subsection{Code converters: Gray code}

The forward conversion is one line: each bit is the XOR of two neighbouring
binary bits. The inverse needs a chain: each next bit depends on the one
already obtained.

\vfile{\codepath/b4_coders/bin2gray.v}{code/b4\_coders/bin2gray.v}

\vfile{\codepath/b4_coders/gray2bin.v}{code/b4\_coders/gray2bin.v}

\textbf{How to read it.} \code{b \textasciicircum\ (b >> 1)} shifts the number
one bit right and XORs it with the original: $g_2=b_2$, $g_1=b_2\oplus b_1$,
$g_0=b_1\oplus b_0$. In the second module the output \code{b[1]} depends on
\code{b[2]}, that is, the chain stretches from the top bit. The testbench
\code{tb\_gray.v} checks that a double conversion returns the number
and that neighbouring codes (including $7\to0$) differ in exactly one bit.

\textbf{Scenario.} The \emph{Waveforms} tab: $b$ from a counter $0$--$7$, the
output \code{g} changes in exactly one bit at every step.

\subsection{Typical mistakes in Verilog}

\textbf{A latch from an incomplete if.} If in a combinational block the
output is left unassigned under some conditions, the circuit ``remembers'' the
previous value; that is, a latch appears: a combinational circuit suddenly
gets memory.

\begin{vcode}
// MISTAKE: for a = 2 and a = 3 nothing is assigned to y
module dec_latch (
    input      [1:0] a,
    output reg [3:0] y
);
    always @* begin
        if (a == 2'd0)      y = 4'b0001;
        else if (a == 2'd1) y = 4'b0010;
    end
endmodule
\end{vcode}

On the \emph{RTL} tab, instead of clean combinational logic a latch
block will appear, and the log will contain a warning about it. The cure: assign a
value \emph{under any condition}: add a final \code{else} or write a
default value at the very beginning of the block.

\begin{vcode}
// FIXED: a default value at the start of the block
module dec_fixed (
    input      [1:0] a,
    output reg [3:0] y
);
    always @* begin
        y = 4'b0000;                       // by default all outputs are zero
        if (a == 2'd0)      y = 4'b0001;
        else if (a == 2'd1) y = 4'b0010;
    end
endmodule
\end{vcode}

\textbf{No default in case.} The same trouble with \code{case}: if not all
values of the selector are listed and there is no \code{default}, the output
is remembered for the rest. If all values are listed (as in $4$ rows for a
$2$-bit \code{a}), \code{default} is not required, but it is still useful to write it.

\begin{vcode}
// MISTAKE: for x other than the four one-hot codes, y is remembered
module enc_latch (
    input      [3:0] x,
    output reg [1:0] y
);
    always @* begin
        case (x)
            4'b0001: y = 2'd0;
            4'b0010: y = 2'd1;
            4'b0100: y = 2'd2;
            4'b1000: y = 2'd3;
        endcase
    end
endmodule
\end{vcode}

\textbf{Blocking versus non-blocking.} In a combinational \code{always @*}
block use \code{=}. The operator \code{<=} is for flip-flops
(Chapter~\ref{ch:ff}); in combinational logic it works, but it
confuses the reader and leads to strange results when the output is used in
the same block.

\begin{diybox}[Complete and change a module]
\textbf{Task.} \stars{1} \emph{Complete it.} Describe an $8\to3$ priority
encoder with an output $v$ using \code{casez}: add eight rows following the
pattern of \code{prio4\_casez}. \stars{2} \emph{Change it.} Add the
hexadecimal digits $A\ldots F$ to \code{seg7} (lower-case \code{b} and
\code{d}, so they are not mistaken for $8$ and $0$). \stars{3} \emph{A stage of the
``Security Panel'' project.} Write a module \code{zone\_display}: four zone
sensors $z$ on the input, seven segments on the output, and for $z=0000$ the
display is blank.

\textbf{In Logra.} The \emph{Code} tab: type your module, press
\emph{Synthesize}. The \emph{Waveforms} tab: for \stars{1} drive $x$ with a
counter $0$--$255$ (or a few typical values: \code{00000000},
\code{00000001}, \code{00110000}, \code{10000000}); for \stars{2} $x$ with a
counter $0$--$15$; for \stars{3} the six rows of the check table.

\textbf{Expected result.} \stars{1} $x=00110000$ gives $y=5$, $v=1$;
$x=0$ gives $v=0$. \stars{2} For $A$ you get \code{1110111}. \stars{3}
For zones 1, 2, 3: \code{0000110}, \code{1011011}, \code{1001111}; for two
zones the higher one is shown; for $z=0$ all zeros.

\textbf{How to check.} For known answers compare the waveform with the
expectations above; the folder \code{code/b4\_coders} contains the testbench
\code{tb\_samples.v}, which checks the model answers (it prints \code{OK}).
\end{diybox}

\textbf{Model answers to the task ``Complete and change a module''.}

\textbf{Model answer to \stars{1}.}
\vfile{\codepath/b4_coders/prio8_casez.v}{code/b4\_coders/prio8\_casez.v}

\textbf{Model answer to \stars{2}} (the other table rows are the same as in
\code{seg7}; here is the whole file):
\vfile{\codepath/b4_coders/seg7_hex.v}{code/b4\_coders/seg7\_hex.v}

\textbf{Model answer to \stars{3}.}
\vfile{\codepath/b4_coders/zone_display.v}{code/b4\_coders/zone\_display.v}

\section{Running project: the zone number on the display}

Let us put everything we have learned into one security panel circuit. There are four zone sensors
$z_3z_2z_1z_0$ (zone~0 means ``nothing triggered''). A priority encoder
turns them into a number ($2$ bits). The display needs a digit in the form of
four bits, so we prepend two zeros to the number on the left. The decoder for the
display converts the digit into seven segments.

\begin{figure}[H]
\centering
\begin{tikzpicture}[node distance=15mm, block/.append style={minimum width=22mm, font=\sffamily\footnotesize}]
  \node[block] (enc) {Priority\\encoder};
  \node[block, right=of enc] (cat) {Concatenation\\\code{00}\,\&\,number};
  \node[block, right=of cat] (seg) {BCD$\to$7\\segments};
  \node[block, right=of seg] (ind) {Display\\$a\ldots g$};
  \draw[arr] ($(enc.west)+(-0.9,0)$) node[left] {$z$ (4)} -- (enc.west);
  \draw[arr] (enc.east) -- node[above, lbl] {2 bits} (cat.west);
  \draw[arr] (cat.east) -- node[above, lbl] {4 bits} (seg.west);
  \draw[arr] (seg.east) -- node[above, lbl] {7 bits} (ind.west);
\end{tikzpicture}
\caption{Stage~4a of the ``Security Panel'' project: from four zone sensors to a digit on the display}
\label{fig:zone_display}
\end{figure}

What the result should be (the higher zone has priority):

\begin{table}[H]
\centering
\caption{Test table for the zone number}
\begin{tabular}{c|c|c}
\toprule
Sensors $z_3z_2z_1z_0$ & Number on the display & Comment \\
\midrule
\code{0000} & 0 & all quiet \\
\code{0010} & 1 & zone 1 triggered \\
\code{0100} & 2 & zone 2 triggered \\
\code{1000} & 3 & zone 3 triggered \\
\code{0110} & 2 & two zones: show the higher one \\
\code{1010} & 3 & two zones: show the higher one \\
\bottomrule
\end{tabular}
\end{table}

\begin{diybox}[Security Panel: stage 4a, the zone number]
\textbf{Task.} \stars{1} Build the chain ``priority encoder
$\to$ concatenation $\to$ BCD$\to$7 segments'' and check all six
rows of the table. \stars{2} Use the encoder output $V$ so that for
$z=0000$ the display goes dark instead of showing ``0'' (pass each segment
through an AND with $V$). \stars{3} \emph{A stage of the ``Security Panel'' project:} instead of
going dark, show a dash ``$-$'' when all is quiet (only segment
$g$ is lit): pass the segments $a\ldots f$ through an AND with $V$, and segment $g$
through an OR with $\overline V$.

\textbf{In Logra.} The \emph{Designer} tab. The blocks ``Priority
encoder'' (width 4), ``Bus concatenation,'' ``Constant,'' and ``BCD\,$\to$\,7
segments'' from the ``Select'' group. Models: the lessons \emph{``Priority
encoder''} and \emph{``A digit on seven segments''}; the example
\code{segment-decoder} (the example library).

\textbf{Expected result.} A seven-bit output: for zone~1 the code
\code{0000110}, for zone~2 \code{1011011}, for zone~3
\code{1001111}; for zone~0 \code{0111111} (in variant \stars{2}
all zeros, in variant \stars{3} \code{1000000}).

\textbf{How to check.} Acceptance criterion: all six rows of the table give the
expected codes on the \emph{Waveforms} tab. On the \emph{RTL} tab you can see
three main blocks. The result of this stage will come in handy in the chapter where
we teach the circuit to compare numbers.
\end{diybox}

\begin{warnbox}[Common mistakes]
\begin{itemize}
  \item \textbf{Mixing up the order of bits.} In a decoder, $a_0$ is the least significant
        bit of the number. If you swap $a_0$ and $a_1$, the wrong line will light, but the
        circuit will still be ``working.'' Test with the number 1 ($01$), not 0 or 3.
  \item \textbf{Forgetting the enable input.} When $E=0$ all outputs are zero.
        If the circuit is ``silent,'' check $E$ first.
  \item \textbf{A simple encoder where several inputs may be active.}
        Two ones on the input give a wrong number. If simultaneous events
        are possible, you need a priority encoder, and $V$ for the ``nothing there'' case.
  \item \textbf{The display shows a ``negative.''} Common anode: a segment
        lights at zero, and you are applying one. You need inverters on all seven
        outputs.
  \item \textbf{An incomplete \code{if} or \code{case} in Verilog.} If for some
        input the output is left unassigned, you get a latch. Assign a default
        value at the beginning of the \code{always @*} block, or write
        \code{else} and \code{default}.
  \item \textbf{\code{<=} in a combinational block.} For combinational
        logic write \code{=}; keep \code{<=} for flip-flops.
\end{itemize}
\end{warnbox}

\begin{ideabox}
\textbf{Decoder}: a number $\to$ a one on one of $2^n$ outputs.
\textbf{Encoder}: a one on one of the inputs $\to$ its number.
\textbf{A priority encoder}, when several inputs are active, chooses the highest.
\textbf{Gray code} changes one bit per step. A digital display is also a
decoder, only its output is not a ``one hot line'' but a set of segments.
\end{ideabox}

\section{Questions and problems}
\label{sec:tasks-coders}

The answers are collected in Appendix~\ref{app:answers}.

\begin{taskbox}[Problems for the chapter]
\begin{enumerate}
  \item \stars{1}~A $3\to 8$ decoder: how many inputs and outputs does it have? Which line is active if the input is 101?
  \item \stars{1}~Write the expressions for all four outputs of a $2\to 4$ decoder with an enable input EN.
  \item \stars{1}~Which segments of a seven-segment display must be lit for the digit 7 (segments $a$--$g$ in the standard layout)?
  \item \stars{1}~A $4\to 2$ priority encoder has inputs $I_3I_2I_1I_0 = 0110$. What code is on the output and why?
  \item \stars{1}~Write the number 6 in binary code and in Gray code. Why is Gray code convenient for angle sensors?
  \item \stars{2}~How many $2\to 4$ decoders are needed to build a $4\to 16$ decoder?
  \item \stars{2}~Find the minimal form of segment $c$ of a seven-segment display, treating the codes $10$--$15$ as don't-care. How many gates are needed?
  \item \stars{2}~Convert the Gray code \code{110} to binary code. Then check the result with the inverse conversion.
  \item \stars{2}~\emph{Find the error.} A student built a $4\to2$ encoder from two OR gates ($y_0=x_1+x_3$, $y_1=x_2+x_3$) and got the code $11$ for the inputs $x_3x_2x_1x_0=0110$. He believes it should be $10$. Who is right? What needs to be changed in the circuit?
  \item \stars{2}~\emph{Find the error.} For the number 5, a student wrote the Gray code \code{100}, ``just by shifting the number.'' What is the actual Gray code of the number 5? Which number does \code{100} correspond to?
  \item \stars{3}~\emph{A project with an acceptance criterion (a stage of the ``Security Panel'' project, 4a).} In the Designer, build the chain ``priority encoder $\to$ concatenation $\to$ BCD$\to$7 segments'' for four zone sensors. Acceptance criterion: all rows of the test table (zone zero, zones 1, 2, 3, two zones at once) give the correct segment code on the Waveforms tab.
  \item \stars{3}~\emph{A project with an acceptance criterion.} Build a $3\to8$ decoder from two $2\to4$ decoders with an enable input. Acceptance criterion: for each of the eight numbers exactly one line is active, and for $E=0$ all outputs are zero.
  \item \stars{3}~\emph{A project with an acceptance criterion.} Build a converter from $3$-bit binary code to Gray code and back. Acceptance criterion: with the double conversion the output equals the input on all eight values, and on the output of the forward converter, while counting $0$--$7$, exactly one bit changes on each step.
  \item \stars{2}~\emph{Find the error (Verilog).} A student described a priority encoder as the chain \code{if (x[3]) y = 2'd3; else if (x[2]) y = 2'd2; else if (x[1]) y = 2'd1; else if (x[0]) y = 2'd0;} inside \code{always @*}, without a final \code{else}. After synthesis a latch block appeared on the RTL tab. Why, and how to fix it?
  \item \stars{3}~\emph{A project with a criterion.} Describe a $3\to8$ decoder with an enable input in Verilog in two ways: in one line with a shift, and from two $2\to4$ decoders (as in Figure~\ref{fig:dc38}). Acceptance criterion: the two modules agree on all 16 combinations $(en,a)$ on the Waveforms tab, and for $en=0$ the outputs are zero.
\end{enumerate}
\end{taskbox}

\section{What's next}

We can now turn a number into a ``one-hot line'' and back, show
a digit on a display, and safely read the position of a sensor. One
related task remains: not just to light one line, but to \emph{select} one of
several signals and pass it to a single output. Today's
decoder, combined with AND gates, turns exactly into such a
selector: a multiplexer. The next chapter is devoted to it.
